\scanpage{115}
\mspage{10\deleted{\unsure{8}}9}
% archive sheet number: 114

Après cette longue \unsure{digression}, il faudrait revenir à l'axiome des supports dans un atelier.

\uline{Lemme} 1 Soit $\Sigma$ un \unsure{qtopos|topos}, $e_{\unsure{i|\ast|\bullet}}=(e_i)_{i\in I}$ une famille d'objets de $\Sigma$ mutuellement disjoints, \deleted{\unsure{Soit telle que}} non nuls, et telle que
\[
\deleted{\sum}\bigvee_i e_i=1_{\Sigma}.
\]

Conditions équivalentes

(i) \unsure{L'ens|L'un} des $e_i$ est contenu dans un sous ensemble $\deleted{\unsure{E}}\unsure{\mathbb{B}|B}_{\unsure{e|\alpha}}$ de $\Sigma$ stable pr \deleted{$\vee$, $\wedge$, $\unsure{\complement|C}$, et qui} $\operatorname{Sup}$, \deleted{et} $\operatorname{Inf}$ et $\complement$, et qui, pour la structure induite et l'unité, est une ``alg. de Boole''.

(ii) \deleted{L'ensemble} Pour tout $J\in\deleted{\unsure{I}}\inserted{\unsure{\mathscr{P}|\mathcal{P}|P}(I)}$, posons
\[
e_J=\operatorname{Sup}_{i\in J}\unsure{e_j|g|e_i}
\]
l'application ainsi obtenue
\[
\mathscr{P}(I)\unsure{\hookrightarrow|\longrightarrow}\Sigma\qquad J\longmapsto e_J
\]
% margin note written slanted
\margin{\uline{NB} Cette application est tjrs injective, car $J=\{i\in I\mid e_i\leq e_J\}$}
(qui commute \unsure{visiblement} aux $\operatorname{Sup}$ quelconques) commute à $\complement$ (donc aussi aux $\operatorname{Inf}$ quelconques).
